\documentclass[11pt,a4paper]{article}

\usepackage[margin=2cm]{geometry}
\usepackage[matrix,frame,arrow]{xypic}
\usepackage{enumitem} %with options [nosep] or [noitemsep]  after \begin{enumerate} etc.
\usepackage[usenames, dvipsnames]{xcolor}
\usepackage{amsmath}
\usepackage{amssymb}
\usepackage{amsthm}
\usepackage{mathrsfs}
\usepackage[alphabetic]{amsrefs}
\usepackage[normalem]{ulem}  %does not alter \emph;   \sout{strike out}; \xout{cross out}
\usepackage[braket,qm]{qcircuit}
\usepackage{caption}  %use \captionsetup{justification=centering} within figure environment
\usepackage{appendix}
\usepackage{yfonts}  %use \textgoth
\def\initdefault{yinit}
\usepackage{tikz}
%\usetikzlibrary{quantikz}  %see http://mirror.aarnet.edu.au/pub/CTAN/graphics/pgf/contrib/quantikz/quantikz.pdf
\usetikzlibrary{matrix,arrows.meta}
%\usepackage{relsize}
%\renewcommand\RSlargest{100pt}
\usepackage{anyfontsize}  %use   \fontsize {⟨size⟩} {⟨baselineskip⟩}
\usepackage{titlesec}
%\titleformat*{\chapter}{\headingfont{16}{18}\selectfont\bfseries}
%\usepackage{CJKutf8}
%\usepackage{ctex}
\usepackage{minitoc}
\usepackage[colorlinks=true, citecolor=ForestGreen, urlcolor=blue]{hyperref} %\href{http://www.ctan.org}{CTAN} produces  output “CTAN”

%\newtheorem{name}[counter]{text}[section]
\theoremstyle{plain} 
\newtheorem{postulate}{Postulate} 
\newtheorem{innercustomthm}{Postulate}
\newtheorem*{prop}{Proposition}
\newenvironment{customthm}[1]
  {\renewcommand\theinnercustomthm{#1}\innercustomthm}
  {\endinnercustomthm}

\theoremstyle{definition}   
\newtheorem*{defn}{Definition} 

\DeclareMathOperator{\Tr}{tr}
\DeclareMathOperator{\Span}{Span}
\renewcommand{\Re}{\operatorname{Re}}
\renewcommand{\Im}{\operatorname{Im}}

\author{John Hutchinson}
\date{\today}
\title{Comments on Leaving Certificate Examinations}

\begin{document}

\maketitle

Until 1966 there were 5 years of high school in New South Wales, the standard age of completion was 17 years, and the final set of (public) exams were known as the \emph{Leaving Certificate Examinations}. It was not uncommon to complete the Leaving Certificate at the age of 16.  Most students, however, left school at the Intermediate Certificate level, after 3 years of high school, and the standard age of completion was then 15.

From 1967, students completed  an additional year of high school, so the standard age of completion was raised to 18.  The new system was known as the \emph{Wyndham Scheme}.  Students typically studied a broader range of subjects than in pre-Wyndham years, but arguably the mathematics/physics/chemistry subjects  were not studied in any greater depth despite the additional year of schooling.

For some comments on the situation for mathematics, see \href{https://www2.merga.net.au/documents/RP162006.pdf}{Leaving Certificate in New South Wales} from 1939 to 1962.

\medskip
Below are the 1962 exam papers for Mathematics I and II Pass Level (which essentially corresponds to ``Mathematical Methods'' in the Australian curriculum) and for Mathematics I and II Honours Level (which essentially correspond to  ``Specialist Mathematics'').

Solutions and comments for the Mathematics II Honours  paper are also included, and I will later do the same for the   Mathematics I Honours paper.

\medskip
It is normally the case that exams are much harder if you havn't actually done the relevant course and lots of similar problems.  So don't be discouraged by looking at the following! 

\medskip


\href{https://maths-people.anu.edu.au/~john/Assets/LC\%20Maths\%20I\%20Pass\%20Exam.pdf}{Mathematics I Pass Level 1962, original paper}

\href{https://maths-people.anu.edu.au/~john/Assets/LC\%20Maths\%20II\%20Pass\%20Exam.pdf}{Mathematics II Pass Level 1962, original paper}

\href{https://maths-people.anu.edu.au/~john/Assets/LC\%20Maths\%20I\%20Honours\%20Exam.pdf}{Mathematics I Honours Level 1962, original paper}

\href{https://maths-people.anu.edu.au/~john/Assets/LC\%20Maths\%20II\%20Honours\%20Exam.pdf}{Mathematics II Honours Level 1962, original paper}

\medskip
\href{https://maths-people.anu.edu.au/~john/Assets/LC_Maths_II_Solns.pdf}{Mathematics II Honours Level 1962, questions, solutions, comments}

 








 

























\end{document} 